% =====================================================================
%  R. Nielsen, PHILOSOPHY AND SOUND HUMAN UNDERSTANDING
%  In: For Idea and Reality.  A Journal, edited by R. Nielsen, B. Bjørnson and Rud. Schmidt.  Year 1870. First Volume (Bind 2).
%  Copenhagen: Chr. Steen & Søn, Publishers; Printed by Louis Klein.  Bind 2, pp. 385-404.  (Filosofi og sund Menneskeforstand, in For Ide og Virkelighed.)
%  ENGLISH TRANSLATION, from transcription.tex (Danish), the text of record.
%  GENERATED by ebook-method/for-idee-1870-73/tr_build.py --write from this template and the translators' parts in
%  ebook-method/for-idee-1870-73/tr/ (TRANSLATOR-BRIEF.md; terms: GLOSSARY.md, its key terms approved by Hans
%  2026-10-13/14).  Conventions: TRANSLATION-PLAYBOOK.md.  Page numbers of the print in the margin.
% =====================================================================
\documentclass[12pt,a4paper,openany]{book}
\usepackage[utf8]{inputenc}
\usepackage[T1]{fontenc}
\usepackage{libertinus}
\usepackage{libertinust1math}
\usepackage{textalpha}
\usepackage[english]{babel}
\usepackage[protrusion=true,expansion=true]{microtype}
\usepackage{setspace}
\usepackage[margin=1.3in]{geometry}
\usepackage{fancyhdr}
\usepackage{hyperref}
\hypersetup{pdftitle={Philosophy and Sound Human Understanding}, pdfauthor={Rasmus Nielsen}, hidelinks}
\setlength{\headheight}{15pt}
\pagestyle{fancy}
\fancyhf{}
\fancyhead[LE,RO]{\thepage}
\fancyhead[C]{\textit{R. Nielsen: Philosophy and Sound Human Understanding (1870)}}
\setstretch{1.3}
\usepackage{marginnote}
\renewcommand{\thefootnote}{\arabic{footnote})}   % the print numbers each page's notes 1), 2)
\newcommand{\opage}[1]{\marginnote{\footnotesize\textit{[#1]}}}
\newcommand{\essayhead}[2]{\begin{center}{\large #1}\ifx&#2&\else\\[0.3em]{\normalsize #2}\fi\\[0.5em]\rule{3em}{0.4pt}\end{center}\medskip\noindent}
\newcommand{\signature}[1]{\par\smallskip\hfill #1\hspace*{2em}\par}

\begin{document}

% the title page of Bind 2 (1870-1-for-idee-og-virkelighed.pdf), in English
\begin{titlepage}
\begin{center}
\vspace*{3em}
{\Huge For Idea and Reality.}\\[1em]
\rule{8em}{0.4pt}\\[1em]
{\large A Journal,}\\[0.8em]
{\small edited}\\[0.4em]
{\footnotesize by}\\[0.6em]
{\large R. Nielsen,}\\[0.4em]
{\large B. Bjørnson and Rud. Schmidt.}\\[1em]
\rule{8em}{0.4pt}\\[0.6em]
{\large Year 1870. First Volume.}\\[0.6em]
\rule{8em}{0.4pt}\\[4em]
Copenhagen.\\[0.4em]
Chr. Steen \& Søn, Publishers.\\[1.5em]
{\small Printed by Louis Klein.}\\[3em]
{\small Translated from the Danish}
\end{center}
\end{titlepage}

\essayhead{Philosophy and Sound Human Understanding.}{}
% --- p. 385 ---
\setcounter{footnote}{0}%
\opage{385}It is still only a little over thirty years since J.
L. Heiberg found himself called upon in our literature to
speak up for a more thoroughgoing philosophical thinking
and to remind us emphatically how far we in this
respect stood behind other civilized nations,
since philosophy was still only a foreign guest in this country, scarcely
even a welcome one. ``It has'',
he says, ``already often been remarked that
while we Danes appropriate the general European
culture in its many different directions, there is yet
one left, by which we have so far not let
ourselves be led, namely the philosophical. The political, the scientific,
the aesthetic and the religious life bear
among us the general stamp of the present; from all those quarters
the beneficent wind of culture blows upon our little Denmark,
which does not merely follow the prevailing
current of air, but also many a time uses sail and
rudder to modify its course, to make use of the
mighty element without letting itself be mastered by it.
But while we willingly let ourselves be blown through by the \emph{spirit}
which, according to the Nordic fable, comes from the four dwarfs
% --- p. 386 ---
\setcounter{footnote}{0}%
\opage{386}in the corners of the world out of Ymir's enormous skull, we
are on the other hand afraid of catching cold from the universal \emph{Spirit},
and cautiously creep into our snail-shell whenever we detect
its presence'' (Perseus. No. 1. 1837).

Much water has run into the shore since that time.
The constitution of the realm, the thoughts of the people, the appearance of literature,
the direction of the spiritual life, everything has in the interval of a
few years undergone such a change that the old
and the new already stand like two disputants each insisting on being right, who watch
one another suspiciously, measure one another with their eyes,
and even where there is good will on both sides,
have difficulty enough in understanding one another. Under the
strong movements in the outward the gaze has at the same time been
turned toward the inward; a growing need for spiritual
self-reflection manifests itself, though in different
ways, in ever wider circles. Thus philosophy
among us is no longer a superfluous luxury, an imported
ornamental plant that can only thrive in a hothouse heated for the purpose,
but an unavoidable necessity, a demand of the spirit,
that we must in our circumstances see to satisfying
as well as we can.

For what is philosophy? A scientific doctrine of consciousness,
a science that wants to clear up what it is
that we call our self, and to help the thinking human being
to orient himself in his own consciousness. All
kinds of delusions and illusions of the senses, disappointments and self-deceptions
of every sort have, humanly speaking, their root
in ignorance of the fundamental laws of the self-conscious life. To
ride astray in one's own forest is dangerous, though in a certain
sense lordly, insofar as it presupposes that
the forest, and with it the estate, must have a wide
% --- p. 387 ---
\setcounter{footnote}{0}%
\opage{387}extent; but to run astray in one's own consciousness and
to ask like the peasant: ``was it I who was here before?'' is
lamentable, and the more so in proportion as the individual
natural gifts are considerable and the spiritual possessions,
knowledge, experiences, etc., extensive. There is therefore,
as one sees, something in philosophy that concerns everyone.
No one can think about his life without philosophizing; reality,
the daily life with its circumstances and
unforeseen accidents, has problems enough: Jeppe in
the Baron's bed is surely close to being choked by problems.

Against the national endeavors of the present time
people have now and then wished to set up the old proposition that truth
is the same for all, that national differences must
not exert any influence on the sciences, that philosophy
in particular, the universal science, must be cosmopolitan.
Insofar as what is still thought of here is the pure laws
of thinking, the proposition certainly has its validity.
But when the question is about the peculiar treatment
of scientific tasks and about the presentation of the content of science
as conditioned by the vernaculars,
nationality asserts its right. The more original
scientific productions are, and the
more the investigations approach what is innermost in the spiritual,
the more strongly will the theory of cognition, and with it the whole
literature, be stamped by the national spirit. That philosophy,
the cosmopolitan science, can by no means be
served by being torn up from the national soil and
floating freely in the airy regions of universal thoughts, its
own history teaches. The philosophy of the Greeks is a reflection of
the Greek spirit, and it is by no means accidental that
% --- p. 388 ---
\setcounter{footnote}{0}%
\opage{388}philosophy at present has a quite different character in
Germany than, e.g., in England.

If then it is true, as is so often repeated,
that it is the clear sound understanding, united with a certain
profundity, but at the same time with a decided aversion to mystical
speculations and idle brooding, that is to
furnish the measure of the intellectual endowment of the Danish people
--- a way of looking at things which finds confirmation in the
fact that Holberg, ``the founder of our literature,
he whose incalculable influence on our taste and our
natural judgment has made itself felt in almost
all directions'', is precisely one of the most
brilliant representatives of the Sound Human Understanding --- if this, we say,
is the case: then there must also be demonstrable
a philosophical task corresponding especially to this direction of spirit.

The task is: \emph{to determine the limits of human}
knowledge.

The first decisive condition for attaining scientific
insight is that at every step I must distinguish precisely
between what I know and what I do not
know. My progress in knowledge depends on this: that there
is much that I do not yet know, but which I may
perhaps in time come to know. In human
knowledge there are on all sides countless limits that
can be overstepped. By new discoveries, new points of view,
new methods, science makes ceaseless progress,
oversteps the old limits in all directions,
and in this sense it will be impossible to predict
how far it may reach in the course of the ages. But the question,
the decisive
% --- p. 389 ---
\setcounter{footnote}{0}%
\opage{389}question of limits for life and science, on which we here in the name of the Sound Human Understanding
must lay stress, is quite another.
The fundamental question is whether philosophy in the form of knowledge is able
to solve the riddles of existence.

Among the riddles of life there are chiefly two, on % print: Livggaaderne (as printed)
whose significance the Sound Human Understanding has always
laid particular stress, and for whose solution it would certainly
be infinitely obliged to philosophy as a science;
these are, as is well known, the reality of God's existence
and the immortality of the soul, for which it seeks confirmation.
Has philosophy been in a position to give such
a confirmation? No! Why not? Because the
last, decisive answer to the questions belonging here
lies outside the limit of human knowledge.
That the many proofs set up hitherto, by which people
have sought to ground those truths, comforting to humanity,
can be dissolved into just as many doubts,
negative criticism has sufficiently demonstrated. By
attacking the proofs of the supposed truths one
has then gone to the opposite extreme; one has
accordingly, from the standpoint of modern consciousness,
set up proofs before whose ``inexorable logic'' both
God and immortality must vanish. God, the almighty
willing self, does not exist; an almighty, willing being
is a conceptless concept, an imagination of faith, a
soap bubble that bursts when it is touched by the needle-point of criticism:
God's existence is, like the existence of the
round square, impossible. In a similar way one proves
that the immortality of the soul is impossible. The life of the soul
is originally conditioned by bodily organs; to think
without a brain is as impossible as to feel without nerves.
% --- p. 390 ---
\setcounter{footnote}{0}%
\opage{390}Thus: when the nerves die, the feeling dies too, when the
body dies, the soul dies too.

But what does the Sound Human Understanding say? It
says that philosophy is presumptuous when it gives itself
the appearance of knowing what no human being can know.
The sound understanding takes sides just as little with the
knowledge that deposes as with that which installs God;
just as little with that which disproves as with that which
proves the immortality of the soul. The sound understanding is,
as regards the highest truths, a great skeptic;
with its grounds for and against it never gets beyond
doubt. But this insuperable doubting of the understanding
contains precisely a presupposition which the understanding itself
must regard as raised above every doubt, the presupposition
that there is in reality something that we in this present
life shall never get to know, that human
knowledge, in other words, has an insuperable limit.

There are many propositions that look reasonable
as long as they are kept in generalities; and yet, when
something definite is spoken of, and it is a question of a real
application, it may suddenly turn out that the propositions
acceptable at first glance contain a hidden
ambiguity and give a false meaning. To this belong
the general propositions of abstract philosophy about the limits
of human knowledge. If there were, one will say,
something altogether inaccessible to knowledge, then there would
have to be, in our own inner life as well as in the world of things,
something altogether inaccessible to thought, something
so enigmatic that it withdrew from all our comprehension, because
it withdrew from all the laws of nature as well as of thought.
But such a mysterious, in itself unthinkable
% --- p. 391 ---
\setcounter{footnote}{0}%
\opage{391}being is a veritable non-being. How dangerous it is to surrender
oneself to unclear representations of a mystical being
whose inexplicable activity is not bound to any
thinkable law, whose revelations mock all concepts
of the connection of things by conditioned causes and
corresponding effects, can be seen from the incessant
struggle of reason with superstition\footnote[1]{The essence of superstition is to assume causes without laws.}. In China, e.g., where a rational
conception of visible things and their connection
has after all been attained to the point that for everyday needs one is tolerably
at home in astronomy, the almanac expressly shows
what a Chinese must suffer every day throughout the year
from secret dread of the enigmatic play of chance; wise
prescriptions are necessary. A Chinese calendar, beginning
with the 8th of February, the Chinese New Year's Day, prescribes:
``On the First you may make your religious offerings
(of fowls as well as fish); you may address
yourself to Heaven (in thanksgivings, prayers, vows --- by
burning gilt paper, straw figures or fireworks in endless
variety); you may dress in gala, with
a fur cap or an elegant waist-band\dots{}. On the Fifth you must
not begin a journey\dots{}. The Eighth is doubtful\dots{}.
On the Tenth you may enter into marriage\dots{}.
On the Thirteenth you must at five o'clock in the morning sit
with your face turned to the southeast''. Without failing to recognize the
great superiority of European culture over the Chinese, one can
nevertheless, when one considers properly how dangerous it still is
in our enlightened age to be thirteen at table, to lay
two knives crosswise, to pick up a pin on a Monday,
to begin something on a Friday, etc., with
% --- p. 392 ---
\setcounter{footnote}{0}%
\opage{392}justice maintain that in the midst of the civilized world
there are still not so very few who, with a shy glance at
the door of chance, which always stands ajar, harbor
a Chinese dread of the whimsical interference of secret powers
in the connection of things.

If the race is ever to get so far that it finally
gets the better of all this superstitious fear and
is rid of all these anxieties, these uncanny forebodings:
then it must gradually learn to free itself
from mystical representations of insoluble riddles. But
what does this mean, other than that the light of thought in
consciousness must at last become so clear that all things
in the world show themselves as necessary revelations of the
eternal reason. The reason in the world and the reason in man
are one and the same being; man's knowledge of
things corresponds exactly to his knowledge of himself. If
it is unthinkable that anything can exist which, by its
nature, falls outside the limit of reason, then it is given thereby
that there cannot possibly be anything either on earth or in heaven
or in the abysses which by its nature falls outside the limit of human knowledge.

With the view here set up all goes excellently
as long as one, as has been said, moves within
the circumference of general considerations. By incessantly
appealing to reason and the laws, the unchangeable laws
of thought and of things, one can even give
the matter the appearance that the strict, investigating sciences,
physics, chemistry, physiology, natural history, etc.,  % print: o, s. v. (as printed)
were all uniting with philosophy to confirm the
principle that human knowledge has no
last, final limit. But if on the other hand we take the matter as
% --- p. 393 ---
\setcounter{footnote}{0}%
\opage{393}it is in reality, then something else shows itself.
For then it shows itself that the principle of the absolute limit
of human knowledge by no means stands in
conflict with the principle of the unchangeableness of the laws;
then it shows itself that the natural sciences, far from
denying the sharply marked limit of knowledge, on the contrary,
all for one and one for all, presuppose it and maintain it
in force.

If we reflect on how it actually came about
that the natural sciences in more recent times have been able to
take so mighty an upswing and attain such great remarkable
results, the ground of this is well enough known,
for it lies in the method. The method, the right procedure,
is determined precisely by this, that one makes clear to oneself
right from the beginning what can be investigated under given
conditions and what lies outside
investigation, what can be known and what cannot
be known. As long as the boundary between what can be
known and what cannot be known was still fluctuating
and indeterminate, as in antiquity and throughout the whole
Middle Ages, so long also the course of the investigations
was wavering, and the yield uncertain; astrology,
alchemy, brooding over the philosopher's stone, such mystical
and magical endeavors show plainly enough how
unclear consciousness was about the limits and
tasks of science.

It is truly remarkable to see the sharpness
with which the strict sciences separate that which
lies outside knowledge from that which can and must become
the object of scientific investigation. ``The aim
of mechanics'', it is said, ``is to fathom the
% --- p. 394 ---
\setcounter{footnote}{0}%
\opage{394}general laws of the equilibrium of the mechanical forces of nature
and of the motions of the masses, whereby a perfect
order and harmony is maintained in the phenomena of the corporeal world''.
One would now think that the hidden essence of the mechanical
forces would have to be cleared up from the ground up
in analytical mechanics; but here the express
declaration is given \emph{that we do not know what force is.}
``When a body'' --- with these words Laplace begins
his ``Mechanics of the Heavens'' --- ``changes its position in relation to
a system of bodies which we regard as at rest,
then it seems to us to move''; and he adds: ``the
nature of the peculiar property which makes a body
pass from one place to another is, and will always
remain, unknown; it has been designated by the name
force, but one can only determine its effects and
the laws according to which it acts''. A similar indication
of the limit of knowledge we meet in chemical physics, e.g.
in the doctrine of heat. ``The sense of feeling is able, in bodies,
to distinguish between different states, which
are expressed by the words: \emph{hot}, \emph{warm}, \emph{cold}, etc.
When a body which one calls cold becomes warm
or hot, its volume grows, it expands. The
\emph{unknown cause,} which brings about the expansion of bodies,
and which at the same time produces the various impressions
mentioned on our feeling, one calls heat''.
Thus: we do not know what heat is; neither do we know
what magnetism is, we do not know what electricity
is; we do not know what the cause of a chemical
attraction is; in short: there is laid down in the foundation of all
doctrinal edifices of natural science an acknowledgment
that human knowledge has an insuperable
% --- p. 395 ---
\setcounter{footnote}{0}%
\opage{395}limit. Herewith the exact sciences have admitted without circumlocution
that existence, despite the natural connection of things, has its riddles; and yet there is hardly
anyone who will maintain that the presupposed limit of knowledge
will instill in human beings a secret dread of the obscure powers of nature, or bring any living soul
to prepare himself for unknown possibilities by sitting
``on the Thirteenth at five o'clock in the morning with his face
turned to the southeast''.

If we consider the relation between soul and body, then
no one can doubt that it is determined in all respects
by laws; none the less the investigations must
begin from two altogether different points of departure.
If the point of departure is chosen in the soul, we get
psychology; if it is chosen in the body, we get physiology.
Psychology keeps to the inner experiences,
to our immediate consciousness of feeling, imagination,
memory, understanding, impulse, will, etc. Physiology
keeps to the outer experiences, it regards
the human body like any other body,
examines its organs and tissues, its inner
structure and composition, partly from a mechanical, partly
from a chemical, partly from a morphological standpoint.

Now it is indeed the case that the investigation of the inward
life of the soul constantly presupposes regard to the
bodily, and that conversely the interpretation of the bodily
formations is conditioned by regard to the functions of the soul;
how would it be possible for the psychologist
to form any concept of the activity of sight
without knowledge of the eye, or for the physiologist to
interpret the structure of the eye without regard to the activity of sight?
% --- p. 396 ---
\setcounter{footnote}{0}%
\opage{396}Psychology and physiology must come to meet each other;
they must diligently confer and correspond with one another.

Nevertheless the two sciences are of so dissimilar a
character that the one forms a part of philosophy,
the other of physics. There is no direct
transition from the one to the other: the transition occurs
by a leap. Feeling and sensory nerve are most intimately
united; but from the physiological analysis of
a sensory nerve to the psychological cognition of a feeling,
what an immense distance! The object of outward
experience is absolutely dissimilar to the object
of inward experience. However long physiology and psychology
continue their investigations, and however much they
may in the course of time approach one another, there
will never come the time when they fully understand one another.
Each speaks its own language, and the two languages are
so immeasurably different that what is said in the one
cannot by any translation be rendered in
the other. It is as though the law and the riddle had here
made an appointment to meet each other; the law can be known, but the riddle
is insoluble. When a sensory nerve is cut, the
part of the body concerned loses all feeling; this is a
law which physiology has discovered, and of which it can
make a rich application; but thereby it has not come
a hair's breadth nearer to answering the question how the nerve
transmits the feeling. And what
likeness is there then between convolutions of the brain, organs of the brain,
the white and the gray matter, and the life of consciousness
of thoughts and representations which, under the supervision of a governing
self, works in the brain?

% --- p. 397 ---
\setcounter{footnote}{0}%
\opage{397}It is altogether impossible either that a brain activity
should be able to produce a thought, or that a thought activity
should cause a movement in the brain, if
there is not an inner common nature in which the two are united.
Such a nature must be bodily like the brain
and spiritual like the thought; but such a nature is
neither an object of thinking nor of sense-perception;
it withdraws from every possible investigation: it lies outside the limit of human knowledge.

And yet there must be such a nature; we know
it with certainty: we can prove it psychologically. If
in the brain activity there were nothing but what we can
imagine by means of physiological observation:
then a brain activity could not possibly produce
a thought. The physiological activity is, whether
it is apprehended from the mechanical or chemical or
morphological side, bodily, solely and only bodily. That
an activity in a bodily part can produce
changes in another bodily part, we can well understand,
for here the cause is of the same kind as the effect;
but that an activity in the merely bodily can produce
a change in the altogether non-bodily is a
contradiction, for here the cause is absolutely dissimilar to
the effect. How one thought can produce
another thought, the thought of ground, e.g., the thought of
consequence, the thought of purpose the thought of means, is obvious,
for here the producing is of the same kind as that
which is produced; but that thought, a non-bodily cause,
should be able to produce in the brain a bodily effect,
is unthinkable, for the cause is again dissimilar,
absolutely dissimilar to the effect.

% --- p. 398 ---
\setcounter{footnote}{0}%
\opage{398}But, one objects, if we know that the connection
between soul and body must absolutely rest on
a nature common to both, a fundamental nature in which the
dissimilar is similar, then we know that such a
nature exists, and can therefore not without contradiction maintain
that its being falls outside our knowledge. The
sophistry of such an objection is obvious. When
we say that there is here something that lies outside our
knowledge, we say quite correctly that we have a knowledge
that such a something is, but we also say that
we cannot get beyond this general knowledge and
determine the unknown more closely; herewith the limit is
fixed.

The soul's law-governed connection with the body is
thus a riddle of life. Simple as this proposition
is, it has nevertheless, as the Sound Human Understanding
will readily see, extremely important consequences. The consequences show
themselves at once when the talk is of the separation of soul
and body, of death; for it goes without saying:
as the connection is, so must the separation be too. The relation
between the law and the riddle becomes visible again.

It is a law, as decisive for psychology
as for physiology, that when the organs of thinking are disturbed,
thinking too is disturbed, and that the light of
consciousness is extinguished when the body dies. Appealing
to this law the naturalists have regarded it
as settled that the soul must of necessity perish
with the body. That is a very hasty assertion.
Everyday experience teaches that consciousness can be
obscured, as e.g. in sleep, indeed even vanish, as
% --- p. 399 ---
\setcounter{footnote}{0}%
\opage{399}in a deep faint, and yet return, because the
life of the soul does not cease. The question whether the soul
is in a position to survive the body is connected most
closely with the question whether the life of the soul,
despite its dependence on the bodily organs, is not
already in this existence raised above perishability
by the fact that in its deeper foundation, in
the hidden nature on which the union of soul and
body originally rests, it has powers and conditions
for forming for itself a new body, a body of a higher
order, when the earthly one is dissolved.

``Man dies like a dog'', says a witty
physiologist; but, we ask, how then does the dog die?
Like a man! That is a fine explanation. The physiologists
themselves admit that they do not know what life is;
we take them at their word, and remind them that they do not
know either what death is. It is thus possible, for here is
the limit of our knowledge, that the soul dies with the body,
but it is also --- for the limit of knowledge cannot
be overstepped --- possible that the soul survives the body.
The objection that a future life cannot be
thought, since every attempt to form a distinct
representation of a higher world-order dissolves into
sheer contradictions, has no validity in strict
science. Physics, which keeps strictly to the natural
facts, constantly runs up against phenomena
whose inner connection it cannot comprehend. It
can, after having found out the laws, lay the acknowledged
laws at the ground of a corresponding explanation,
and yet the explanation may dissolve into
% --- p. 400 ---
\setcounter{footnote}{0}%
\opage{400}contradiction\footnote[1]{An example. Ampère explains the magnet as consisting of manifold molecular, parallel and similarly directed electric circuits. From this hypothesis all magnetic phenomena can be interpreted in a manner exactly corresponding to the mathematical formulas. But --- these circuits are now supposed to take place in the midst of a conductor (the steel, e.g.), and must then, in accordance with Ampère's own law, leave the molecule and spread over the whole body. The state can therefore not continue, unless the molecule constantly produces a new circuit. But --- an electric circuit presupposes a corresponding consumption and a corresponding production of heat; but in the magnet there is no production of heat and no consumption, therefore also no electric circuit.}. Were natural science to deny the phenomena whose
essence it cannot comprehend, it would have to end by
denying the whole of nature.

On the secure determination of the limit of knowledge
rests the relation between faith and knowledge. When the outermost
limit of human knowledge is reached, the solution
of the highest tasks of life will always present itself in such a way
that there appear two possibilities, equally valid and yet canceling
one another: it is possible that the soul
is mortal, and it is possible that it is immortal, it is
possible that there is a living, personal, willing God,
but the opposite is also possible. Only when science,
after all-sided testing, has given grounded
certainty that the one of the two possibilities has the same
scientific validity as the other, has knowledge said
its last word and must leave the decision to faith.
Faith now chooses the one of the two possibilities, unbelief
the other. The separation between faith and unbelief
takes place by the choice; the choice is not scientific, but
absolutely personal, here the choice is made by virtue of personal
grounds.

% --- p. 401 ---
\setcounter{footnote}{0}%
\opage{401}On the foundation determined by the choice philosophical
doctrinal edifices can again be erected; but these doctrinal edifices
now differ from the directly scientific ones
in that the grounds are grounds of personality, in that it
is not necessity but freedom on which the grounding
rests. The riddle and the law thereby enter
into a new phase.

It is an objection as well known as it is worn out
that the doctrine of personality, by making, in the name of freedom,
faith an independent principle of cognition, introduces arbitrariness,
violates the laws and comes into conflict
with the naturally necessary connection of things. In particular
it is the doctrine of miracle that arouses offense. It
may, one says, from the standpoint of science well be admitted
that there is much that is enigmatic in nature. Natural science
indeed constantly discovers new phenomena, new
forces and new possibilities; but these discoveries can
by no means weaken our conviction of the necessary connection
of things, since every new discovery is, if
not always a discovery of new laws, yet an extension
and confirmation of those already known. With
the principle of will, on the other hand, arbitrariness is introduced; but
arbitrariness, the watchword of wild possibilities, is incompatible
with the laws.

How matters stand with this question can be
illuminated from a human as well as from a divine
side. We ascribe to human beings, and not only to them
but also to animals, arbitrary movements, movements
which have their cause in immediate self-determination.
If it were now really true that every manifestation of
arbitrary self-determination ran counter to the laws of nature, then
% --- p. 402 ---
\setcounter{footnote}{0}%
\opage{402}one would necessarily have to assume one of two things: either arbitrariness
is an appearance, and there is nothing arbitrary at all,
or, if there are arbitrary manifestations of life, then
arbitrariness must at every moment violate the laws.
If both assumptions must be declared false, then it is
thereby also given that the arbitrariness of self-determining beings
may very well be reconciled with unchangeable
laws.

There is in existence a thoroughgoing opposition
between the essence of nature and of spirit, of necessity and
of freedom, but there is in the principle itself just as little
any conflict between the laws of the two regions as
there is within the compass of the realm of nature any conflict between
the mechanical and chemical laws.

Yes, one continues from the side of criticism, Such things can
well be heard in this way in general, and as long as they are only heard
in all generality, it sounds quite good; but if we bring forward
the cases, the absurdity shows itself at once. That
the arbitrariness of the human will cannot break
the connection of nature follows directly from the fact that man's
will is impotent; but according to the supreme fundamental thought
of the doctrine of personality, God's almighty will
is supposed to be able to intervene arbitrarily in the order of nature, to perform
miracles and thus to break the laws. --- Here again
is a mixture of true and false. It is quite certainly a
distinguishing mark of the miracle that it has the almighty will
as its freely determining cause; but it is, from the standpoint of faith and
of personality, by no means the mark of the miracle
that it should involve a breach of the divine
laws. To make the contrary-to-law the mark of the miracle
is an invention of the more recent criticism. The reason
% --- p. 403 ---
\setcounter{footnote}{0}%
\opage{403}why this invention has gained such strong currency that the proposition
is now close to being exalted into one of the great
fundamental principles of science, lies plainly open.
In order to look like science and to botch the explaining
or at least the defending of the miracle, theology has from
time out of mind had to overstep the limit of knowledge; in
order to refute the botched explanation and to be rid of the miracle,
negative criticism has for its part likewise
been tempted to overstep the limit of knowledge. Theology
concludes: if God is almighty, then he must be able
at pleasure to annihilate the natural causes
and suspend the laws; criticism concludes: if God performs
a deed that cannot be explained according to the laws known to us,
then he does a misdeed and comes into
contradiction with himself. That is false logic\footnote[1]{By altering the expression and saying: an event that conflicts with a single one of the laws known to us conflicts with all laws, one wins only an ambiguous advantage. The expression, to conflict with a law, can namely be taken in two senses. It can mean: to stand in contradiction to the concept of law as law (to draw a round square, to exhibit an effect without a cause, to set 2 + 2 = 5 and the like). In this sense the expression finds no application to the miracle. It can next mean: to be in apparent conflict with a real law, so that the fulfillment of the unknown law seems to cancel the known one. In this sense the expression finds application not only to the miracle, but also to natural events among themselves. Let us imagine that someone who had hitherto learned only to explain phenomena by mechanical laws suddenly became a witness of surprising phenomena of magnetic, electric, electromagnetic laws: would he not then think that he saw nothing but wonders?}. The
empty reasoning about what God can do and what
God cannot do confuses the brains.

According to what has here been adduced, anyone with sound
understanding and ordinary cultivation will now be able to say to himself
% --- p. 404 ---
\setcounter{footnote}{0}%
\opage{404}how important it is that once for all a
limit be set between what can be known and what cannot
be known; only when the limit is set, and its validity
acknowledged, will life and science be able to meet
in honest understanding and look one another in the eye. A task
of such thoroughgoing significance will naturally
gradually press forward and be treated in all
countries by all thinkers, as far as science
reaches; but that does not exclude that its solution may
lie nearer to the one national individuality than to the other.
German speculation has gone beyond the limit;
English empirical philosophy has not reached the limit;
French logic of the understanding, with its exact and yet
mobile forms, makes the limit fluid and abstract by
making it mathematical; it is then not unreasonable that
a people whose peculiar literature has a Holberg
for its father and a profound mythology for its mother, a people
with whose plain sound understanding it so naturally accords
that it has not hurried itself into building heaven-storming
systems, but has given itself time and begun only
to philosophize on its own account after it has secured in natural science,
history and poetry a sound
foundation for the doctrine of spirit, that, we say, such a people
might perhaps have its eyes opened to one of the most essential
tasks of the present age, that its late-born philosophy might grow
with the task and further its solution by contributing, if
only a contribution; a contribution is always something.
\signature{R. \emph{Nielsen}.}

\end{document}
